% Part: proof-theory
% Chapter: natural-deduction
% Section: introduction

\documentclass[../../../include/open-logic-section]{subfiles}

\begin{document}

\olfileid{pt}{nat}{int}

\olsection{جان پچھان}

قدرتی اخذ !!{proof} نظاماں دا اک ٹبّر اے جیہدے وچ ہر منطقی جوڑ یا مسوّر لئی استنباطاں دی اک جوڑی اے۔ اک استنباط عامل نوں «شامل» کردا اے (!!a{introduction} قاعدہ) تے دوجا عامل نوں «ہٹاؤندا» اے (!!a{elimination} قاعدہ)۔ مثلاً \Intro{\lif} دا نتیجہ !!{formula}s ہوندے نیں~$!A \lif !B$ دی شکل دے، تے \Elim{\lif} دے مقدمیّاں وچ !!{formula}s~$!A \lif !B$ دی شکل دے ہوندے نیں۔

قدرتی اخذ دے پہلے دو روپ 1930 دی دہائی وچ گیرہارٹ گینٹسن تے ستانیسواف یاشکوفسکی (Stanisław Jaśkowski) نے پیش کیتے۔ ثبوت دے نظریہ دان زیادہ تر گینٹسن دے نظام دی گل کردے نیں، تے یاشکوفسکی دے نظام توں اوہ نظام نکلے جیہڑے پڑھائی وچ سبھ توں ودھ ورتے جاندے نیں۔ دوہاں نظاماں وچ \emph{مفروضیاں} توں شروع کر سکدے آں۔ ایہ اوہ !!{formula}s نیں جیہناں نوں !!{prove}d کرن دی لوڑ نہیں، پر اوہ !!{proving} وچ، ہور !!{formula}s دے، ورتے جا سکدے نیں۔ پورا !!{proof} ایہو جیہے مفروضیاں اُتے ٹِک سکدا اے۔ چونکہ اوہ ثابت نہیں کیتے، !!{proof} اپنا نتیجہ (اپنا \emph{آخری !!{formula}}) قائم نہیں کردا، بس ایہ قائم کردا اے کہ نتیجہ مفروضیاں توں نکلدا اے۔

کجھ استنباطی قاعدیاں دے نتیجے کجھ مفروضیاں اُتے ہور نہیں ٹِکدے۔ ایہ قاعدے مفروضے «خارج» کردے نیں۔ مثلاً \Intro{\lif} قاعدہ !!a{proof} لیندا اے نتیجے~$!B$ دا، مفروضے~$!A$ توں، تے ایہدا نتیجہ $!A \lif !B$ اے۔ !!{proof} جیہڑا $!A \lif !B$ اُتے مُکدا اے،~$!A$ اُتے ہور نہیں ٹِکدا۔ مفروضہ $!A$ خارج ہو جاندا اے۔ گینٹسن دے نظاماں وچ !!{proof}s، !!{formula}s دے درخت نیں، مفروضے سبھ توں اُتلے !!{formula}s نیں، تے~\Intro{\lif} ورگے قاعدیاں اُتے نشان نیں جیہڑے دسدے نیں کہ کیہڑے مفروضے خارج ہوندے نیں۔ یاشکوفسکی دے نظاماں وچ !!{proof} دے اوہ حصے جیہڑے کسے مفروضے اُتے ٹِکدے نیں، کسے ڈھنگ نال وکھ رکھے جاندے نیں، مثلاً کھڑی لکیر پچھے ہٹواں لکھ کے یا خانے وچ پا کے۔ خارج کیتا مفروضہ~$!A$ تے شرطیے دی تالی $!B$ خانے دیاں پہلی تے آخری سطراں نیں، تے نتیجہ~$!A \lif !B$، جیہڑا \Intro{\lif} دا اے، خانے توں باہر اے۔

نظام «قدرتی» نیں کیوں جے استنباطی قاعدے ساڈی (یا گھٹ توں گھٹ ریاضی داناں دی) قدرتی دلیل ورگے نیں۔ فرضی نتیجہ (شرطیہ) قائم کرن لئی مقدم منّ کے اوس نال تالی ثابت کردے آں۔ ایہ~\Intro{\lif} اے۔ جدوں شرطیہ $!A \lif !B$ تے اوہدا مقدم~$!A$ معلوم نیں، تاں نتیجہ~$!B$ کڈھدے آں۔ ایہ~\Elim{\lif} اے۔ فصلی عبارت~$!A \lor !B$ منّ کے کوئی نتیجہ سچ وکھاؤن لئی صورتاں وکھ کر کے ثبوت دیندے آں۔ ایہ~\Elim{\lor} اے۔ عام دعویٰ~$\lforall[x][!A(x)]$ ثابت کرن لئی وکھاؤنے آں کہ $!A(c)$ سچ اے، کسے من مرضی دے~$c$ لئی۔ ایہ~\Intro{\lforall} اے۔ ایس طرح اگے۔

مثلاً گینٹسن طرز دے قدرتی اخذ وچ سوکھا !!{proof}~$!A \lif (!A \lor !B)$ دا ایہ اے:
\[
\AxiomC{$\Discharge{!A}{x}$}
\RightLabel{\Intro{\lor}}
\UnaryInfC{$!A \lor !B$}
\DischargeRule{\Intro{\lif}}{x}
\UnaryInfC{$!A \lif (!A \lor !B)$}
\DisplayProof
\]
ایہ مفروضے~$!A$ توں شروع ہوندا اے، \Intro{\lor} لا کے~$!A
\lor !B$ کڈھدا اے، تے فیر \Intro{\lif} لا کے~$!A \lif (!A \lor !B)$ کڈھدا اے۔ \Intro{\lif} استنباط مفروضہ~$!A$ خارج کردا اے۔ ایہ نشان~$x$ دَسدا اے۔

ثبوتی نظریے دے ماہر گینٹسن دے اصلی نظام پسند کردے نیں جیہڑے !!{formula}s دے درختاں نال کم کردے نیں، بھاویں اک وکھرا روپ وی ثبوتی نظریے وچ اہم اے۔ !!^a{proof}~$!A$ دا، مفروضیاں~$\Gamma$ توں، وکھاؤندا اے کہ $!A$،~$\Gamma$ توں نکلدا اے، یعنی $\Gamma \Entails !A$۔ قدرتی اخذ دے استنباطی قاعدے قدرتی طور اُتے ایہو جیہے منطقی نتیجیاں بارے گل کردے پڑھے جا سکدے نیں۔ مثلاً \Intro{\lif} کہندا اے کہ جے $\Gamma + !A
\Entails !B$، تاں $\Gamma \Entails !B$۔ قاعدیاں نوں سدھا ایویں لکھنا قدرتی اے کہ اوہ استنباطی قاعدے دے مقدمیّاں تے نتیجے وچ، مفروضیاں تے اوہناں توں نکلن والے نتیجے دوہاں اُتے کم کرن۔ یعنی اوہ سیکوئنٹاں $\Gamma \Sequent !A$ اُتے کم کرن۔ اُتلا سوکھا ثبوت ایہو جیہے نظام وچ ایویں دِسے گا:
\[
\Axiom$x:!A \fCenter !A$
\RightLabel{\Intro{\lor}}
\UnaryInf$x:!A \fCenter !A \lor !B$
\DischargeRule{\Intro{\lif}}{x}
\UnaryInf$\fCenter !A \lif (!A \lor !B)$
\DisplayProof
\]
جیویں تسی ویکھدے او، قاعدے اوہی نیں۔ اوہ !!{formula} اُتے کم کردے نیں~$\Sequent$ دے سجّے۔ کھبّے !!{formula}s بس اوہناں مفروضیاں دی فہرست نیں جیہناں اُتے ہر قدم ٹِکدا اے۔

!!{introduction}/!!{elimination} قاعدیاں دیاں ساریاں جوڑیاں دا سیٹ اپنے آپ وچ کلاسیکی منطق لئی پورا نہیں۔ سانوں دو ہور قاعدے شامل کرنے پین گے جیہڑے~$\lfalse$ نال جڑے نیں۔ پہلا «وجدانی محال دا قاعدہ»~\FalseInt، یا «دھماکا»، کہندا اے کہ $\lfalse$ توں کجھ وی کڈھ سکدے آں۔ دوجا بالواسطہ ثبوت دا کلاسیکی اصول~\FalseCl{} (یا «کلاسیکی محال دا قاعدہ») کہندا اے کہ جے اسیں !!{prove} کریئے~$\lfalse$ نوں~$\lnot !A$ توں، تاں اسیں !!{prove}d کر لَیا اے~$!A$۔ جے~\FalseCl چھڈ دیئے، تاں وجدانی منطق لئی ڈھکدا نظام ملدا اے۔ جے دوویں چھڈ دیئے، تاں «کمترین منطق» ملدی اے۔

اسیں کلاسیکی تے وجدانی منطق لئی گینٹسن طرز دے قدرتی اخذ دی گل کراں گے۔ ایہ نظام \Log{N1c} تے \Log{N1i} نیں۔ نظام \Log{N2c} تے \Log{N2i} اوہناں نال میل کھاندے نیں، پر سیکوئنٹاں $\Gamma \Sequent !A$ اُتے کم کردے نیں، اکّلے !!{formula}s~$!A$ دی تھاں۔

\end{document}
